\section{Opis zagadnienia}
\label{sec:zagadnienie}

\subsection{Sformułowanie problemu}
\label{sec:problem}

Środowisko modelowane jest jako siatka, gdzie każde pole reprezentuje
fragment terenu o~określonym koszcie przejścia.

\begin{itemize}
    \item Każde pole posiada wagę określającą trudność terenu.
    \item Koszt przejścia zależy od~kierunku ruchu:
    \begin{itemize}
        \item ruch prosty (poziomy/pionowy): mnożnik $1$,
        \item ruch po~skosie (diagonalny): mnożnik $1{,}4$.
    \end{itemize}
\end{itemize}

Na planszy rozmieszczone są \textbf{atrakcje turystyczne}:
\begin{itemize}
    \item każda atrakcja ma wartość turystyczną,
    \item każda atrakcja ma koszt odwiedzenia (np.\ cena biletu),
    \item każda atrakcja należy do~określonego \emph{typu} (zabytek,
          przyroda, restauracja itp.).
\end{itemize}

\textbf{Podróżnik:}
\begin{itemize}
    \item startuje w~punkcie $s$,
    \item kończy w~punkcie $t$,
    \item posiada ograniczony czas $T$ oraz budżet $B$.
\end{itemize}

\textbf{Dodatkowe wymagania konstrukcyjne:}
\begin{itemize}
    \item każda atrakcja może być odwiedzona maksymalnie raz,
    \item dla każdego typu atrakcji zdefiniowane są limity dolny $l_{\tau}$
          i~górny $u_{\tau}$ na liczbę atrakcji tego typu w~trasie.
\end{itemize}

Problem stanowi rozszerzenie klasycznego zagadnienia znajdowania ścieżki
o~elementy decyzyjne związane z~wyborem punktów pośrednich --- bliski
\emph{Orienteering Problem} z~ograniczeniami typowymi.

\subsection{Model matematyczny}
\label{sec:model}

\subsubsection{Definicja grafu}

Problem reprezentowany jest jako graf:
\[
G = (V, E)
\]
gdzie:
\begin{itemize}
    \item $V$ --- zbiór pól siatki,
    \item $E$ --- możliwe przejścia (krawędzie) między polami sąsiadującymi
          w~siatce 8-kierunkowej (sąsiedztwo Moore'a).
\end{itemize}

\subsubsection{Koszt przejścia}

Koszt przejścia pomiędzy polami:
\[
c_{ij} = w_i \cdot d_{ij}
\]
gdzie:
\begin{itemize}
    \item $w_i$ --- waga pola $i$,
    \item $d_{ij} \in \{1,\ 1{,}4\}$ --- współczynnik kierunku ruchu
          ($1$ dla ruchu osiowego, $1{,}4$ dla diagonalnego).
\end{itemize}

\subsubsection{Atrakcje}

Niech:
\begin{itemize}
    \item $A \subseteq V$ --- zbiór atrakcji,
    \item $v_i$ --- wartość atrakcji $i$,
    \item $p_i$ --- koszt (cena) atrakcji $i$,
    \item $\tau_i \in \mathcal{T}$ --- typ atrakcji $i$.
\end{itemize}

\subsubsection{Typy atrakcji}

Niech $\mathcal{T}$ oznacza zbiór wszystkich typów atrakcji. Każdemu
typowi $\tau \in \mathcal{T}$ przypisane są następujące parametry:
\begin{itemize}
    \item $t_{\tau}$ --- czas potrzebny na zwiedzenie atrakcji typu $\tau$,
    \item $l_{\tau}$ --- minimalna liczba atrakcji typu $\tau$, które
          należy odwiedzić,
    \item $u_{\tau}$ --- maksymalna liczba atrakcji typu $\tau$, które
          można odwiedzić.
\end{itemize}

\subsubsection{Zmienne decyzyjne}

\[
x_{ij} =
\begin{cases}
1 & \text{jeśli przechodzimy z } i \text{ do } j, \\
0 & \text{w przeciwnym razie,}
\end{cases}
\qquad
y_i =
\begin{cases}
1 & \text{jeśli odwiedzamy atrakcję } i, \\
0 & \text{w przeciwnym razie.}
\end{cases}
\]

\subsubsection{Funkcja celu}

Celem jest maksymalizacja wartości odwiedzonych atrakcji:
\begin{equation}
\max \;\; \sum_{i \in A} v_i \cdot y_i .
\label{eq:obj}
\end{equation}

\subsubsection{Ograniczenia}

\textbf{1.\ Start:}
\begin{equation}
\sum_{j} x_{s j} = 1 .
\label{eq:start}
\end{equation}

\textbf{2.\ Koniec:}
\begin{equation}
\sum_{i} x_{i t} = 1 .
\label{eq:end}
\end{equation}

\textbf{3.\ Zachowanie przepływu:}
\begin{equation}
\sum_{i} x_{i k} = \sum_{j} x_{k j}
\quad \forall k \in V \setminus \{s, t\} .
\label{eq:flow}
\end{equation}

\textbf{4.\ Ograniczenie czasu:}
\begin{equation}
\sum_{(i,j) \in E} c_{ij} \cdot x_{ij}
+ \sum_{i \in A} y_i \cdot t_{\tau_i}
\;\le\; T .
\label{eq:time}
\end{equation}

\textbf{5.\ Ograniczenie budżetu:}
\begin{equation}
\sum_{i \in A} p_i \cdot y_i \;\le\; B .
\label{eq:budget}
\end{equation}

\textbf{6.\ Ograniczenia ilościowe na typy atrakcji:}
\begin{equation}
l_{\tau} \;\le\; \sum_{\substack{i \in A \\ \tau_i = \tau}} y_i \;\le\; u_{\tau}
\quad \forall \tau \in \mathcal{T} .
\label{eq:types}
\end{equation}

\textbf{7.\ Powiązanie atrakcji ze~ścieżką.}
Atrakcja może być odwiedzona tylko wtedy, gdy ścieżka przez nią przechodzi:
\begin{equation}
y_i \;\le\; \sum_{j} x_{j i} \quad \forall i \in A .
\label{eq:link}
\end{equation}

\textbf{8.\ Jednokrotność odwiedzin atrakcji:}
\begin{equation}
y_i \le 1 \quad \forall i \in A
\label{eq:attr-unique}
\end{equation}

\textbf{9.\ Binarność:}
\begin{equation}
x_{ij} \in \{0, 1\}, \quad y_i \in \{0, 1\} .
\label{eq:bin}
\end{equation}

Równania \eqref{eq:obj}--\eqref{eq:bin} formalnie definiują problem
optymalizacji jako mieszany model przepływu w~grafie z~elementami doboru
podzbioru (klasy \emph{Orienteering Problem}).

\subsection{Uwagi modelowe i decyzje projektowe}

\begin{itemize}
    \item \textbf{Złożoność.} Problem jest NP-trudny (zawiera w~sobie
          uogólnione \emph{Orienteering Problem}), co uzasadnia
          zastosowanie metod heurystycznych i~meta-heurystycznych
          --- w~tym projekcie algorytmu \emph{Artificial Bee Colony}
          (patrz~rozdział \ref{sec:algorytmy}).
    \item \textbf{Ścieżka jako konstrukcja pomocnicza.} W~implementacji
          fragmenty trasy pomiędzy waypointami (start, kolejne atrakcje,
          koniec) generowane są algorytmem A\textsuperscript{*}
          (najtańsza ścieżka, koszt zgodny z~$c_{ij}$, sąsiedztwo
          8-kierunkowe). Algorytm uwzględnia dynamiczne zakazy: blokowane
          są pola atrakcji niezaplanowanych w~bieżącej liście waypointów,
          Algorytm A* nie nakłada ograniczeń na ponowne
          wykorzystanie tej samej krawędzi.
    \item \textbf{Atrakcje są blokujące.} W~implementacji pole atrakcji,
          której nie planujemy odwiedzić, jest blokowane dla ścieżki ---
          jeśli komórka atrakcji znajdzie się na trasie, atrakcja zostanie
          ,,zaliczona''. Wymusza to spójność: trasa nie biegnie tranzytem
          przez atrakcje, których nie zamierzaliśmy odwiedzić.
    \item \textbf{Kara vs.\ dopuszczalność.} W~zaproponowanym podejściu nie
          stosujemy funkcji kary za naruszenie ograniczeń: algorytm utrzymuje
          jedynie rozwiązania w~pełni dopuszczalne (spełniające
          \eqref{eq:start}--\eqref{eq:bin}), a~rozwiązania niedopuszczalne
          są odrzucane już na etapie konstrukcji.
    \item \textbf{Kryterium porównawcze.} Do porównywania rozwiązań
          w~populacji ABC stosowana jest leksykograficzna krotka
          \[
              \big(\,
                  \textstyle\sum v_i y_i,\;
                  -(\text{czas trasy}),\;
                  -(\text{koszt ruchu}),\;
                  -(\text{koszt atrakcji})
              \,\big),
          \]
          gdzie minus oznacza preferencję niższej wartości. Pozwala to przy
          jednakowej wartości celu różnicować rozwiązania krótsze czasowo
          i~tańsze, co ułatwia generowanie zróżnicowanej populacji.
\end{itemize}
